% TEMPLATE-MILC-v3.tex - plantilla maestra UNICA de las guias MILC v3.
%
% La usan las tres familias de guias, cada una con su driver:
%   · Grados 6-11  -> scripts/build-guias-g11.py
%   · Bebras       -> scripts/build-guias-bebras.py
%   · Semillero    -> scripts/build-guias-semillero.py
%
% Antes habia tres copias casi identicas (~400 lineas iguales, 6 distintas):
% cada mejora habia que aplicarla tres veces, y en la practica no se hacia
% (el motor de recursos vivia solo en dos de ellas). Lo que varia entre
% familias esta parametrizado en 6 placeholders que cada driver llena
% (se nombran aqui SIN su sintaxis real: el builder reemplaza por texto plano,
% tambien dentro de los comentarios, y un valor multilinea rompe el preambulo):
%
%   HEADER_IZQ        encabezado izquierdo de pagina
%   HEADER_DER        encabezado derecho de pagina
%   PORTADA_ETIQUETA  rotulo sobre el numero grande (GRADO/MOMENTO/MODULO)
%   PORTADA_SERIE     linea de serie de la portada
%   PORTADA_ID        identificador de la guia en la portada
%   PORTADA_CIRCULO   el circulito de grado (°), o vacio si no aplica
%
% El resto de claves las documenta content/guias/_SCHEMA.md. El script valida
% que no queden placeholders sin reemplazar antes de escribir el archivo final.

\documentclass[11pt,letterpaper]{article}
\usepackage[margin=2.0cm]{geometry}
\usepackage[table]{xcolor}
\usepackage{fontspec}
\usepackage[spanish]{babel}
\usepackage{tikz}
% Librerías TikZ para los diagramas de las guías (bloque `recursos:`):
%  · babel  → babel en español vuelve activos caracteres como «>», y TikZ los usa
%             en sus opciones (>=stealth, ->). Sin esta librería, cualquier
%             diagrama con flechas revienta con «\language@active@arg> has an
%             extra }». Debe ir ANTES de que se dibuje nada.
%  · positioning        → sintaxis «below left=2pt and 3pt».
%  · decorations.pathreplacing → llaves ({brace}) para señalar rangos.
\usetikzlibrary{babel,positioning,decorations.pathreplacing}
\usepackage{graphicx}
% Circuitos: el trabajo de electrónica se hace hoy con micro:bit y sus sensores
% integrados (MakeCode, sin cableado), así que no se necesitan esquemáticos.
% Si algún día entran montajes con protoboard, basta descomentar esta línea:
% los diagramas .tex/.tikz declarados en `recursos:` ya se incrustan solos.
% \usepackage{circuitikz}
\usepackage[most]{tcolorbox}
\usepackage{tabularx}
\usepackage{array}
\usepackage{fancyhdr}
\usepackage{titlesec}
\usepackage[document]{ragged2e}  % bandera a la derecha en todo el cuerpo
\usepackage{enumitem}
\usepackage{microtype}

% Helvetica Neue no trae la flecha U+2192, asi que cada «→» escrito literalmente
% en el YAML se compilaba a NADA, en silencio (462 flechas en 61 guias: rutas del
% tipo «paleta → tipografia → lockup» salian rotas). Mapeamos el caracter a su
% version matematica, que si tiene glifo. Asi el autor puede seguir escribiendo
% «→» comodamente en el YAML.
\usepackage{newunicodechar}
\newunicodechar{→}{\ensuremath{\rightarrow}}
% Mismo problema con los simbolos de marca y vinetas: el «✓» de «una palomita ✓»
% salia como cuadrito vacio en el PDF de todas las guias que lo usaban.
\usepackage{pifont}
\newunicodechar{✓}{\ding{51}}
\newunicodechar{✗}{\ding{55}}
\newunicodechar{✔}{\ding{52}}
\newunicodechar{•}{\textbullet}
\newunicodechar{≥}{\ensuremath{\geq}}
\newunicodechar{≤}{\ensuremath{\leq}}

\setmainfont{Helvetica Neue}
\setsansfont{Helvetica Neue}
\setlength{\parindent}{0pt}
\setlength{\parskip}{6pt}
% Sin particion de palabras: en 6.º-7.º una palabra cortada por guion al final
% de linea («Comuni-cacion») frena la lectura mucho mas que una linea corta.
\hyphenpenalty=10000
\exhyphenpenalty=10000
\pretolerance=2000
\tolerance=3000
\setlength{\emergencystretch}{3em}
\linespread{1.10}
\renewcommand{\arraystretch}{1.25}
\setlist{leftmargin=5mm,itemsep=1mm,topsep=1mm}
\newcolumntype{Y}{>{\RaggedRight\arraybackslash}X}

\definecolor{milcMagenta}{HTML}{E6007E}
\definecolor{milcVino}{HTML}{5A0038}
\definecolor{milcTurquesa}{HTML}{0093A5}
\definecolor{milcVerde}{HTML}{58A923}
\definecolor{milcMostaza}{HTML}{E5B400}
\definecolor{milcOcre}{HTML}{B86600}
\definecolor{milcNegro}{HTML}{241816}
\definecolor{milcGris}{HTML}{F7F7F4}
\definecolor{milcLinea}{HTML}{E8E2DD}
\definecolor{milcGradePrimary}{HTML}{0066FF}
\definecolor{milcGradeDark}{HTML}{003D99}
\definecolor{milcGradeSoft}{HTML}{D6E8FF}
\definecolor{milcGradeAccent}{HTML}{CFE7FF}
\definecolor{milcGradeText}{HTML}{FFFFFF}
\color{milcNegro}

\pagestyle{fancy}
\fancyhf{}
\lhead{\small Tecnología e Informática · Grado 11}
\rhead{\small Guía 20 · MILC v3}
\cfoot{\small\thepage}
\renewcommand{\headrulewidth}{0pt}

\newtcolorbox{titlebox}[1]{
  enhanced,colback=#1,colframe=#1,arc=7mm,boxrule=0pt,
  left=7mm,right=7mm,top=3mm,bottom=3mm,
  before skip=8pt,after skip=8pt,
  fontupper=\sffamily\bfseries\Large\color{white}
}

\newtcolorbox{softbox}[3]{
  enhanced,colback=#2,colframe=#1,borderline west={4pt}{0pt}{#1},
  arc=2mm,boxrule=.4pt,left=4mm,right=4mm,top=3mm,bottom=3mm,
  before skip=6pt,after skip=8pt,title={#3},coltitle=white,
  colbacktitle=#1,fonttitle=\sffamily\bfseries,fontupper=\RaggedRight,
  attach boxed title to top left={xshift=3mm,yshift=-2mm},
  boxed title style={arc=2mm, boxrule=0pt}
}

\newtcolorbox{drawbox}[2]{
  enhanced,colback=white,colframe=#1,arc=4mm,boxrule=1pt,
  left=5mm,right=5mm,top=4mm,bottom=4mm,before skip=8pt,after skip=8pt,
  title={#2},coltitle=white,colbacktitle=#1,fonttitle=\sffamily\bfseries,
  fontupper=\RaggedRight,
  attach boxed title to top left={xshift=4mm,yshift=-2mm},
  boxed title style={arc=3mm, boxrule=0pt}
}

\newtcolorbox{infoband}[1]{
  enhanced,colback=#1!8,colframe=#1!50,arc=2mm,boxrule=.5pt,
  left=4mm,right=4mm,top=2mm,bottom=2mm,before skip=4pt,after skip=4pt,
  fontupper=\small\RaggedRight
}

% Puente narrativo entre secciones (contrato editorial Conéctate).
% Da continuidad: "Acabas de X. Ahora vas a Y porque Z."
% Visualmente: banda izquierda en color del grado, texto en itálica pequeña.
\newtcolorbox{puentebox}{
  enhanced,colback=milcGradeSoft,colframe=milcGradeDark,
  borderline west={3pt}{0pt}{milcGradeDark},
  arc=0mm,boxrule=0pt,left=5mm,right=4mm,top=2.5mm,bottom=2.5mm,
  before skip=4pt,after skip=8pt,
  fontupper=\itshape\small\color{milcNegro}\RaggedRight
}
% La flecha va en modo matematico a proposito: Helvetica Neue NO tiene el glifo
% U+2192, asi que \textrightarrow se compilaba a NADA («Missing character: There
% is no → in font Helvetica Neue») y el puente salia sin flecha. En modo
% matematico la flecha viene de la fuente matematica y si se dibuja.
\newcommand{\puente}[1]{%
  \begin{puentebox}%
    \textcolor{milcGradeDark}{$\rightarrow$}\hspace{2mm}#1%
  \end{puentebox}%
}

\newcommand{\linead}[1]{\noindent\color{milcLinea}\rule{#1}{0.5pt}\color{milcNegro}}
\newcommand{\respuesta}[1]{\par\vspace{2mm}\foreach \n in {1,...,#1}{\noindent\color{milcLinea}\rule{\linewidth}{0.5pt}\par\vspace{5mm}}\color{milcNegro}}
\newcommand{\checkbox}{\textcolor{milcGradeDark}{\fbox{\phantom{x}}}\hspace{5pt}}

\newcommand{\phasecard}[4]{
\begin{tcolorbox}[enhanced,colback=#3,colframe=#2,arc=3mm,boxrule=.5pt,height=3.05cm,valign=center,halign=center,left=1.5mm,right=1.5mm,top=2mm,bottom=2mm]
{\sffamily\bfseries\fontsize{22}{24}\selectfont\textcolor{#2}{#1}}\par
{\sffamily\bfseries\small #4}
\end{tcolorbox}
}

% ── Piezas del rediseno 2026 (legibilidad 6.º-11.º) ───────────────────
% Banda de estacion: reemplaza el titlebox macizo. Titulo a la izquierda,
% dato practico (tiempo/modalidad) a la derecha, en una sola linea.
\newcommand{\milcestacion}[2]{%
\begin{tcolorbox}[enhanced,colback=#1,colframe=#1,arc=2mm,boxrule=0pt,
  left=5mm,right=5mm,top=2.6mm,bottom=2.6mm,before skip=11pt,after skip=7pt]
{\sffamily\bfseries\fontsize{15}{18}\selectfont\color{white}#2}
\end{tcolorbox}}

% Tarjeta de paso numerada: sustituye las tablas «Pilar / Aplicacion», que
% eran formato de consulta docente, no de estudiante. El numero va en un
% circulo sobre el borde para que la secuencia se lea de un vistazo.
\newcommand{\milcpaso}[3]{%
\begin{tcolorbox}[enhanced,colback=white,colframe=milcLinea,arc=1.5mm,boxrule=.9pt,
  left=14mm,right=4mm,top=3mm,bottom=3mm,before skip=4pt,after skip=4pt,
  fontupper=\RaggedRight,
  overlay={\node[anchor=north west,circle,fill=#2,text=white,inner sep=0pt,
    minimum size=7.4mm,font=\sffamily\bfseries\small]
    at ([xshift=3.6mm,yshift=-3.2mm]frame.north west) {#1};}]
#3
\end{tcolorbox}}

% Casilla marcable de tamano util con lapiz (la anterior era \fbox{\phantom{x}},
% de alto variable segun el texto que la seguia).
\newcommand{\milcmarca}{\framebox[3.8mm]{\rule{0pt}{3.4mm}}}

% Fila marcable de lista suelta (checks de la estacion 1).
\newcommand{\milccheckrow}[1]{%
\par\vspace{1.8mm}\noindent
\makebox[6.4mm][l]{\raisebox{-.55mm}{\textcolor{milcNegro}{\milcmarca}}}%
\parbox[t]{\dimexpr\linewidth-6.4mm\relax}{\RaggedRight #1}\par}

% Celda marcable dentro de la rubrica (tabla de autoevaluacion). Misma
% sangria francesa, pero pensada para vivir dentro de una columna X.
\newcommand{\milcopcion}[1]{%
\makebox[5.2mm][l]{\raisebox{-.55mm}{\textcolor{milcNegro}{\milcmarca}}}%
\parbox[t]{\dimexpr\linewidth-5.2mm\relax}{\RaggedRight #1}}

% ── Recursos visuales de la guía ──────────────────────────────────────
% Declarados en el bloque `recursos:` del YAML y emitidos por el builder.
% \guiaFigura   → imagen rasterizada o PDF (foto, carta celeste, montaje).
% \guiaDiagrama → fuente TikZ/circuitikz (.tex) incrustada como vector:
%                 es la vía para circuitos y esquemáticos editables.
% #1 = ruta relativa al .tex · #2 = pie (puede ir vacío)
\newcommand{\guiaFigura}[2]{%
\begin{center}
\includegraphics[width=0.88\linewidth,height=0.38\textheight,keepaspectratio]{#1}\\[1.5mm]
{\footnotesize\itshape\textcolor{milcNegro!75}{#2}}
\end{center}
}

\newcommand{\guiaDiagrama}[2]{%
\begin{center}
\input{#1}\\[1.5mm]
{\footnotesize\itshape\textcolor{milcNegro!75}{#2}}
\end{center}
}

\begin{document}
\thispagestyle{empty}

% ── Portada ───────────────────────────────────────────────────────────
% Antes: un rectangulo de color a pagina completa. Se veia bien en pantalla,
% pero en la fotocopiadora del colegio se comia el toner y salia gris sucio.
% Ahora la banda ocupa el tercio superior y el resto va en blanco.
\begin{tikzpicture}[remember picture,overlay]
  \path[fill=milcGradePrimary]
    (current page.north west) rectangle ([yshift=-10.6cm]current page.north east);

  \node[anchor=north west,text=milcGradeText] at ([xshift=2.0cm,yshift=-1.55cm]current page.north west)
    {\sffamily\bfseries\fontsize{12}{14}\selectfont GRADO};
  \node[anchor=north west,text=milcGradeText,opacity=.85] at ([xshift=2.0cm,yshift=-2.35cm]current page.north west)
    {\sffamily\bfseries\fontsize{8.5}{10}\selectfont SERIE GUÍAS · TECNOLOGÍA E INFORMÁTICA};
  \node[anchor=north east,text=milcGradeText,opacity=.85,align=right] at ([xshift=-2.0cm,yshift=-1.55cm]current page.north east)
    {\sffamily\bfseries\fontsize{9}{12}\selectfont I.E. Sor María Juliana\\Cartago, Valle del Cauca};

  \node[anchor=west,text=milcGradeText] (gradonum) at ([xshift=1.85cm,yshift=-7.1cm]current page.north west)
    {\sffamily\bfseries\fontsize{112}{112}\selectfont 11};
  % El circulito de grado (°) solo aplica a las guias de grado («11°»); en
  % Bebras y Semillero el numero grande es un momento/modulo, no un grado.
  % Se posiciona relativo al nodo `gradonum`, no en coordenadas absolutas.
  \draw[milcGradeText,line width=1.6pt]
    ([xshift=2.0mm,yshift=-4.5mm]gradonum.north east) circle (.15cm);

  \node[anchor=west,text=milcGradeText,text width=11.4cm,align=left] at ([xshift=1.5cm,yshift=0.5cm]gradonum.east)
    {
     {\sffamily\bfseries\fontsize{9}{11}\selectfont\MakeUppercase{Período 2 · Automatización y procesos}}\\[3mm]
     {\sffamily\bfseries\fontsize{21}{25}\selectfont\setlength{\baselineskip}{27pt}Caso real ---\\[4pt]automatizar\\[4pt]un proceso real}\\[3.5mm]
     {\sffamily\bfseries\fontsize{15}{18}\selectfont Guía 20-11-TIC}
    };
\end{tikzpicture}

\vspace*{9.4cm}
\vfill

{\sffamily\bfseries\fontsize{8}{10}\selectfont\textcolor{milcGradeDark}{LO QUE TE LLEVAS HOY}}

\begin{tcolorbox}[enhanced,colback=white,colframe=milcNegro,arc=2mm,boxrule=1.6pt,
  left=5mm,right=5mm,top=4mm,bottom=4mm,before skip=3pt,after skip=12pt,
  fontupper=\RaggedRight\fontsize{11}{15.5}\selectfont]
1 proceso real automatizado de principio a fin (en el colegio, la casa o un microemprendimiento del barrio), con evidencia antes/después documentada (KPIs medidos antes y después, capturas de pantalla del antes vs después, retroalimentación de al menos 2 usuarios reales), bitácora del ciclo completo (mapa → BPMN → Form → hoja → trigger → IA → bot → KPIs → mejora) y video de 2-3 minutos sustentando el caso
\end{tcolorbox}

\begin{tcolorbox}[enhanced,colback=milcGris,colframe=milcGris,arc=2mm,boxrule=0pt,
  left=5mm,right=5mm,top=3.5mm,bottom=3.5mm,after skip=12pt]
{\sffamily\bfseries\fontsize{8}{10}\selectfont\textcolor{milcNegro!55}{ESTA GUÍA ES DE}}\\[3mm]
\begin{tabularx}{\linewidth}{X p{3.2cm} p{3.2cm}}
\linead{\linewidth} & \linead{3.0cm} & \linead{3.0cm}\\[-1mm]
{\fontsize{8.5}{10}\selectfont\textcolor{milcNegro!55}{Nombre completo}} &
{\fontsize{8.5}{10}\selectfont\textcolor{milcNegro!55}{Grupo}} &
{\fontsize{8.5}{10}\selectfont\textcolor{milcNegro!55}{Fecha}}\\
\end{tabularx}
\end{tcolorbox}

\vfill

\noindent\textcolor{milcLinea}{\rule{\linewidth}{1pt}}\\[2mm]
{\sffamily\bfseries\fontsize{9.5}{12}\selectfont Autor: Dr. Álvaro Cárdenas Orozco}\\[1mm]
{\sffamily\fontsize{8.5}{11}\selectfont\textcolor{milcNegro!55}{Metodología MILC · apertura ancestral · pensamiento computacional · triángulo Dussel–estoicismo–Floridi}}

\newpage

\section*{Apertura: Caso real --- automatizar un proceso del colegio o la casa}

\begin{softbox}{milcGradeDark}{milcGradeSoft}{Saber ancestral}
En los pueblos del Valle del Cauca había una figura que hoy se subestima: \textbf{el reparador del barrio}. El \textbf{relojero} que ajustaba el tic-tac de los relojes de pared con paciencia de horas; el \textbf{soldador} que arreglaba el silbato del olla express, el cuadro de la bicicleta, la tapa del fogón; la \textbf{costurera} que recosía el pantalón gastado o le subía el dobladillo. Tenían tres valores en común: (1) \textbf{paciencia} para entender qué falla; (2) \textbf{precisión} para tocar solo lo necesario; (3) \textbf{devolver algo útil} a quien encargó el trabajo. No inventaban máquinas nuevas: reparaban lo que estaba roto en su entorno. Esa tradición sigue viva en el oficio digital: hoy se reparan procesos que \emph{ya casi funcionan}, pero pierden tiempo, dinero o paciencia. Tu capstone retoma ese oficio.
\end{softbox}

\begin{softbox}{milcTurquesa}{milcGris}{Saber contemporáneo}
Un \textbf{caso real de automatización} es la integración de todo lo aprendido en el periodo (mapeo, BPMN, Form, hoja-base, trigger, IA, bot, KPIs, mejora continua) aplicado a un proceso \textbf{que existe, que duele a alguien real, y que está al alcance de tu mano}. No se trata de inventar un emprendimiento ficticio: se trata de ofrecer a un docente, un acudiente, un emprendedor del barrio, o tu propia casa, un proceso \emph{visiblemente mejor} que el actual. La evidencia es la regla del oficio: \textbf{KPIs antes y KPIs después} comparados, capturas del flujo antes y después, y al menos 2 usuarios reales que validen el cambio. El video de sustentación (2-3 min) es el acto público que cierra el periodo --- y entrena la habilidad profesional de defender el propio trabajo.
\end{softbox}

\begin{softbox}{milcMostaza}{milcGris}{Pregunta-puente}
\textit{¿Cómo sabía el reparador del barrio cuándo decir \emph{``ya está''} y cuándo seguir ajustando? ¿Y qué pierde un emprendedor novato cuando entrega un proyecto sin \textbf{evidencia antes/después}, sin usuarios reales que lo validen y sin haber medido lo que cambió?}
\end{softbox}

\begin{softbox}{milcVerde}{milcGris}{Saber y saber hacer}
Al terminar podrás: (1) \textbf{identificar} un proceso real cercano que necesita automatización y describirlo con claridad; (2) \textbf{analizar} el estado actual con KPIs antes de intervenir; (3) \textbf{crear} la automatización completa integrando las 9 sesiones del periodo en una sola implementación coherente; (4) \textbf{evaluar} el resultado con KPIs después + 2 usuarios reales + video de sustentación, declarando con honestidad qué funcionó y qué no.
\end{softbox}



\small
\begin{tabularx}{\textwidth}{>{\bfseries}p{3.0cm}Y}
\rowcolor{milcGradePrimary!12}Autor & Dr. Álvaro Cárdenas Orozco\\
DBA & Diseño y mejora de procesos digitales con criterio ético (MEN, T\&I)\\
Referentes & Pensamiento computacional · Lean / Kaizen · Floridi (ética informacional)\\
Producto final & 1 proceso real automatizado de principio a fin (en el colegio, la casa o un microemprendimiento del barrio), con evidencia antes/después documentada (KPIs medidos antes y después, capturas de pantalla del antes vs después, retroalimentación de al menos 2 usuarios reales), bitácora del ciclo completo (mapa → BPMN → Form → hoja → trigger → IA → bot → KPIs → mejora) y video de 2-3 minutos sustentando el caso\\
\end{tabularx}
\normalsize

\puente{Antes de empezar, mira el viaje completo. Has aprendido 9 piezas a lo largo del periodo: mapeo, BPMN, formularios, hojas-base, triggers, IA, chatbots, KPIs, Kaizen. Hoy las integras en \textbf{un solo caso real} con evidencia. Es el capstone --- y la prueba de que el oficio digital, como el del relojero, se mide por la pieza terminada que funciona en manos del usuario.}

\milcestacion{milcGradeDark}{Ruta MILC: así vamos a aprender}
\begin{tabularx}{\textwidth}{>{\centering\arraybackslash}X>{\centering\arraybackslash}X>{\centering\arraybackslash}X>{\centering\arraybackslash}X}
\phasecard{1}{milcVerde}{milcGris}{Escucha\\[-1mm]\small Observo, escucho y pregunto.} &
\phasecard{2}{milcTurquesa}{milcGris}{Entender\\[-1mm]\small Sistematizo y ordeno.} &
\phasecard{3}{milcMagenta}{milcGris}{Hacer\\[-1mm]\small Construyo y aplico.} &
\phasecard{4}{milcVino}{milcGris}{Cierre\\[-1mm]\small Evalúo y reflexiono.}
\end{tabularx}


\puente{Antes de teorizar, vas a hacer una \emph{auditoría rápida} de procesos reales cercanos para elegir uno: tu casa, el colegio, un microemprendimiento del barrio. Vas a anotar 3 candidatos y a evaluar cuál duele más a alguien real, cuál tiene KPI medible y cuál puedes implementar en una semana.}

\milcestacion{milcVerde}{Estación 1 de 3 · Escucha}

\begin{softbox}{milcVerde}{milcGris}{Escena para escuchar}
\textbf{Actividad 1 · IDENTIFICA --- 3 candidatos a capstone} (15 min · individual). Sin filtrar todavía por viabilidad, anota 3 procesos reales cercanos que duelan: en tu casa (recordatorios de medicación, control de gastos, asignación de tareas), en el colegio (registro de asistencia, encuesta a padres, comunicación con grado), en el barrio (pedidos del tendero, reservas de un peluquero, agendamiento de mototaxis). Para cada candidato anota: \emph{¿a quién duele?, ¿qué KPI se puede medir antes?, ¿puedo implementar en 1 semana?}. Elige uno como capstone con justificación de 1 frase.
\end{softbox}

{\sffamily\bfseries\fontsize{8}{10}\selectfont\textcolor{milcNegro!55}{MARCA CUANDO LO HAYAS HECHO}}
\milccheckrow{Anoté 3 candidatos sin censurarlos por dificultad.}
\milccheckrow{Evalué cada candidato en 3 dimensiones (dolor, KPI, viabilidad).}
\milccheckrow{Elegí uno con justificación de 1 frase.}

\begin{infoband}{milcVerde}


\textbf{Lo que va al cuaderno (Actividad 1 · IDENTIFICA).} Título: «Actividad 1 --- 3 candidatos a capstone». \textbf{Formato:} tabla 3 filas × 4 columnas (Candidato~|~A quién duele~|~KPI medible~|~Implementable en 1 semana). Al final 1 línea con el elegido y por qué. \textbf{Sabes que terminaste cuando} tienes claro \emph{a quién} le entregarás el caso al final del periodo.
\end{infoband}

\puente{Ya tienes candidato. Ahora vas a entender la \textbf{anatomía} del caso real completo: los 7 entregables mínimos del capstone, los errores típicos (caso ficticio, sin medición antes/después, sin usuarios reales, sin video), y la regla de oro de la evidencia profesional.}

\milcestacion{milcTurquesa}{Estación 2 de 3 · Entender}

\begin{softbox}{milcTurquesa}{milcGris}{Lo que vas a entender}
El capstone parece complejo porque integra todo, pero su anatomía es ordenada. Vamos a descomponer los 7 entregables mínimos con los pilares del pensamiento computacional para que tu caso pase la prueba profesional.
\end{softbox}

\milcpaso{1}{milcTurquesa}{El capstone se descompone en 7 entregables: (1) \textbf{Caso elegido}: proceso real + dueño + dolor concreto. (2) \textbf{Mapa BPMN}: estado actual diagramado (S2 del periodo). (3) \textbf{Implementación}: Form, hoja, trigger, IA o bot --- al menos 3 de las 6 piezas integradas. (4) \textbf{KPIs antes/después}: mismos 3-5 KPIs medidos en ambos momentos. (5) \textbf{2 usuarios reales validando}: testimonios escritos o en video. (6) \textbf{Bitácora del proceso}: qué se hizo cada día, qué falló, qué se ajustó. (7) \textbf{Video de sustentación 2-3 min}: caso, antes/después, decisión propuesta.}
\milcpaso{2}{milcTurquesa}{Todo capstone sigue el patrón \textbf{``proceso real $\to$ dolor real $\to$ usuario real $\to$ evidencia real''}: si una pieza es ficticia, todo el capstone es ficticio. La diferencia entre un proyecto académico y un trabajo profesional es la realidad de las 4 piezas. La phronesis del oficio exige resistir la tentación del \emph{``como si''}: \emph{``como si tuviera clientes''}, \emph{``como si midiera''}, \emph{``como si funcionara''}.}
\milcpaso{3}{milcTurquesa}{Lo esencial cabe en una pregunta: \emph{¿el caso le mejoró la vida a alguien real?}. Si la respuesta es no, el capstone está incompleto aunque las 7 piezas técnicas estén. La evidencia humana (los 2 usuarios) es lo que separa el ejercicio del oficio.}
\milcpaso{4}{milcTurquesa}{Algoritmo: (1) elige el caso; (2) mapea el estado actual; (3) mide KPIs antes; (4) implementa al menos 3 piezas; (5) pide a 2 usuarios reales que usen el proceso nuevo; (6) mide KPIs después; (7) escribe la bitácora; (8) graba el video; (9) entrega al dueño del proceso (no solo al docente).}

\begin{drawbox}{milcTurquesa}{Anatomía del capstone bien hecho}
\textbf{Caso:} dueño + dolor + KPI inicial. \\[1mm] \textbf{Mapa BPMN:} estado actual diagramado. \\[1mm] \textbf{Implementación:} 3+ piezas técnicas integradas. \\[1mm] \textbf{Evidencia antes/después:} mismos KPIs comparados. \\[1mm] \textbf{Usuarios reales:} 2 testimonios escritos o filmados. \\[1mm] \textbf{Bitácora:} cronograma + fallos + ajustes. \\[1mm] \textbf{Video:} 2-3 minutos con caso, antes/después y propuesta.
\end{drawbox}

\begin{softbox}{milcMostaza}{milcGris}{Errores comunes y su corrección}
(1) Caso ficticio o inventado: rompe el oficio profesional. (2) KPIs solo después (sin antes): no hay evidencia de mejora. (3) Sin usuarios reales: queda como ejercicio de aula. (4) Bitácora ausente: nadie puede reproducir ni discutir el trabajo. (5) Video que no muestra antes/después: el lector solo ve el resultado, no la transformación. (6) Implementar las 9 piezas mal por ambición: mejor 3 piezas funcionando bien que 9 a medias.
\end{softbox}

\begin{infoband}{milcTurquesa}


\textbf{Lo que va al cuaderno (Actividad 2 · ANALIZA).} Título: «Actividad 2 --- Anatomía del capstone». \textbf{Formato:} ficha con los 7 entregables aplicados a tu caso elegido + 6 errores comunes con marca 🚨 del que más temes + cronograma semanal de implementación. \textbf{Sabes que terminaste cuando} tienes un plan claro de qué entregar y cuándo.
\end{infoband}

\puente{Tienes el método. Toca implementarlo: 9 piezas del periodo integradas en un proceso real, KPIs antes y después medidos, 2 usuarios reales validando, bitácora del ciclo completo, video de 2-3 min. Es trabajo de varios días --- la sesión de hoy abre el camino y deja la planificación.}

\milcestacion{milcMagenta}{Estación 3 de 3 · Hacer}

\begin{softbox}{milcMagenta}{milcGris}{Cómo vas a construir tu producto}
\textbf{Actividad 3 · CREA + EVALÚA --- Plan completo del capstone} (30 min · individual). Hora de planificar. La implementación dura una o dos semanas; hoy dejas el plan con cronograma realista, los KPIs antes ya medidos, la entrevista a los 2 usuarios agendada, y la estructura del video definida.
\end{softbox}

\milcpaso{1}{milcMagenta}{Tu plan se descompone en 5 piezas: (1) caso elegido con dueño/dolor/KPI; (2) cronograma diario de 5-7 días con qué tarea hacer cada día; (3) KPIs antes ya medidos con los datos reales actuales; (4) agenda de las 2 entrevistas con usuarios; (5) guion del video de 2-3 minutos (qué decir en cada minuto).}
\milcpaso{2}{milcMagenta}{Sigue el patrón \textbf{``medir antes ANTES de intervenir''}: nunca cambies nada del proceso actual hasta haber medido los 3-5 KPIs base. Sin ese antes, no podrás demostrar ningún después. Esa disciplina inicial vale por tres horas de trabajo posterior.}
\milcpaso{3}{milcMagenta}{El guion del video se estructura en 3 actos: (1) \textbf{Caso} (30 seg): \emph{``este es el proceso de X, lo sostiene Y, dolía porque Z''}; (2) \textbf{Antes/Después} (60-90 seg): mostrar los KPIs comparados + capturas; (3) \textbf{Propuesta} (30 seg): \emph{``a este caso le sigue Z''}.}
\milcpaso{4}{milcMagenta}{Algoritmo: (1) escribe el caso en 3 líneas; (2) construye cronograma de 5-7 días; (3) mide KPIs antes hoy; (4) agenda usuarios; (5) escribe guion del video; (6) prepara el ambiente de implementación (cuentas Zapier, Sheets, ManyChat, etc.); (7) entrega el plan al docente para retroalimentación temprana.}

\begin{drawbox}{milcMagenta}{Checklist antes de entregar el plan}
\checkbox Caso elegido con dueño/dolor/KPI claro. \\[1mm] \checkbox Cronograma de 5-7 días con tareas concretas por día. \\[1mm] \checkbox KPIs antes medidos hoy con datos reales. \\[1mm] \checkbox 2 usuarios reales agendados para validar después. \\[1mm] \checkbox Guion del video estructurado en 3 actos. \\[1mm] \checkbox Bitácora abierta (cuaderno o hoja) para ir registrando avance. \\[1mm] \checkbox Cuentas y herramientas listas (ManyChat, Zapier, Sheets, etc.).
\end{drawbox}

\begin{softbox}{milcMagenta}{milcGris}{Plantilla de trabajo}
\textbf{Caso.} \textit{Proceso + dueño + dolor + KPI inicial.} \\[1mm] \textbf{Cronograma.} \textit{Día 1: medir antes. Día 2-4: implementar. Día 5: probar con usuarios. Día 6: medir después. Día 7: video.} \\[1mm] \textbf{KPIs antes.} \textit{Lista con valores actuales.} \\[1mm] \textbf{Usuarios.} \textit{Nombre + relación + cuándo entrevistar.} \\[1mm] \textbf{Video.} \textit{Guion en 3 actos.}
\end{softbox}

\begin{infoband}{milcMagenta}


\textbf{Lo que va al cuaderno (Actividad 3 · CREA + EVALÚA).} Título: «Actividad 3 --- Plan del capstone». \textbf{Formato:} ficha con 5 piezas del plan + cronograma + KPIs antes medidos + guion del video. \textbf{Extensión:} 2 páginas de cuaderno. \textbf{Sabes que terminaste cuando} podrías empezar la implementación mañana sin pedir más instrucciones.
\end{infoband}

\puente{Lo que sigue resume \emph{qué} entregas al final del periodo: 1 proceso automatizado con evidencia antes/después + bitácora + video. El criterio principal: que un emprendedor o docente real, viendo el video, pueda decir \emph{``este caso me sirve''} y pedir el contacto del estudiante para conversar.}

\milcestacion{milcMostaza}{Producto de la sesión}
\textbf{Plan completo del capstone + KPIs antes medidos + cronograma}.
\begin{softbox}{milcMostaza}{milcGris}{Criterios de calidad}
(1) Caso real con dueño, dolor y KPI claros. (2) Cronograma de 5-7 días con tareas concretas. (3) KPIs antes ya medidos con datos reales. (4) 2 usuarios agendados. (5) Guion del video en 3 actos. (6) Herramientas y cuentas listas para implementar.
\end{softbox}

\puente{Antes de cerrar, mira el capstone desde las cinco dimensiones humanas. Cerrar bien es respeto por el oficio entero; entregar sin evidencia es desprecio por quienes confiaron en ti.}

\milcestacion{milcVino}{Evaluación liberadora · cinco dimensiones}

\begin{softbox}{milcVino}{milcGris}{Sentido de la evaluación}
Evaluar no es solo revisar una nota. Es reconocer qué aprendí, qué cambió en mí, qué emociones manejé, cómo me relaciono con la ciudad y la comunidad, y qué le dejaré a quien viene detrás.
\end{softbox}

\textbf{Mi reflexión liberadora en cinco dimensiones.} Completa cada frase iniciada:

\small
\begin{tabularx}{\textwidth}{>{\bfseries\RaggedRight\arraybackslash}p{3.4cm}Y>{\RaggedRight\arraybackslash}p{4.6cm}}
\rowcolor{milcVino}\color{white}Dimensión & \color{white}\bfseries Mi respuesta (frame starter) & \color{white}\bfseries Algo que descubrí\\
1. Personal & Lo que más cambió en mí fue \linead{6.5cm} & \linead{4.2cm}\\
2. Emocional & La emoción más fuerte fue \linead{2.5cm} y la regulé \linead{3cm} & \linead{4.2cm}\\
3. Ciudadana & En democracia esto importa porque \linead{6cm} & \linead{4.2cm}\\
4. Local (Cartago) & En mi barrio sería útil para \linead{6.5cm} & \linead{4.2cm}\\
5. Intergeneracional & A mi abuela/abuelo le diría \linead{6.5cm} & \linead{4.2cm}\\
\end{tabularx}
\normalsize

\begin{infoband}{milcVino}
\textbf{Para la próxima sustentación.} Vuelve a esta tabla en seis meses. Verás que las respuestas cambian — no porque lo aprendido cambie, sino porque tú cambias. La evaluación liberadora es una conversación contigo a través del tiempo.
\end{infoband}


\puente{Y para terminar, tres voces sobre el reparador y el capstone. Dussel sobre los procesos que liberan o sobrecargan, Epicteto sobre la honestidad de declarar lo que no funcionó, Floridi sobre la entrega con evidencia como ética profesional contemporánea.}

\milcestacion{milcGradeDark}{Triángulo de pensamiento · cierre filosófico}

\begin{softbox}{milcVino}{milcGris}{Dussel · lente del nosotros}
\textit{Toda automatización decide si libera al trabajador o lo reemplaza, si cuida al usuario o lo descarga.}

{\footnotesize\itshape\color{milcNegro!70}--- formulación de esta guía, al modo de Enrique Dussel. No es una cita textual.}

Tu capstone es una toma de posición ética. Cuando automatizas un proceso, decides quién se beneficia: el dueño que ahorra tiempo, el operador que pierde su tarea, o el usuario que recibe mejor servicio. La phronesis del cierre exige declarar en el video, con honestidad, \emph{a quién benefició el cambio} y \emph{a quién (si a alguien) le agregó trabajo o le quitó algo}. Esa transparencia es la diferencia entre el oficio profesional y el truco técnico.

\textbf{En mi capstone, ¿quién gana y quién pierde con la automatización? ¿Mi video lo declara o lo esconde?}
\end{softbox}
\linead{\linewidth}\\[1mm]
\linead{\linewidth}

\begin{softbox}{milcMostaza}{milcGris}{Estoicismo · Epicteto · cuidado interior}
\textit{La honestidad de reportar lo que no funcionó es la prueba final del oficio.}

{\footnotesize\itshape\color{milcNegro!70}--- formulación de esta guía, al modo de Epicteto. No es una cita textual.}

Es tentador presentar el capstone solo con sus triunfos. La sabiduría estoica recomienda lo contrario: \emph{declara también lo que falló, lo que ajustaste, lo que aún no resolviste}. Esa honestidad es ética profesional. El emprendedor o docente que recibe tu caso valorará mucho más la verdad incompleta que el éxito decorado. La transparencia es virtud técnica.

\textbf{En mi video de sustentación, ¿cuento solo lo que funcionó o también lo que falló, lo que tuve que rehacer, lo que no logré? ¿Quién aprende más del relato honesto?}
\end{softbox}
\linead{\linewidth}\\[1mm]
\linead{\linewidth}

\begin{softbox}{milcTurquesa}{milcGris}{Floridi · lente de la infoesfera}
\textit{La entrega con evidencia es la ética profesional del oficio digital contemporáneo.}

{\footnotesize\itshape\color{milcNegro!70}--- formulación de esta guía, al modo de Luciano Floridi. No es una cita textual.}

Una entrega sin evidencia (sin KPIs antes/después, sin usuarios reales, sin bitácora) es declaración de fe, no oficio profesional. La cultura digital contemporánea exige evidencia reproducible: lo que no se mide ni se muestra, no existe. Cerrar el periodo con un capstone con evidencia te entrena en esa ética que rige toda profesión técnica del XXI.

\textbf{¿Mi capstone tiene la evidencia que un consultor profesional o un inversionista exigirían? ¿O sigue siendo trabajo escolar disfrazado de profesional?}
\end{softbox}
\linead{\linewidth}\\[1mm]
\linead{\linewidth}

\begin{infoband}{milcNegro}
\textbf{Nota para el docente.} Las tres formulaciones son del autor de la guía, no citas textuales de los pensadores: recogen su lente para pensar el tema, no sus palabras. Si vas a citarlos en clase, remite a la obra original.
\end{infoband}

\milcestacion{milcVino}{Mi autoevaluación · marca una casilla por fila}

{\small
\setlength{\extrarowheight}{2.2pt}
\begin{tabularx}{\textwidth}{>{\RaggedRight\arraybackslash\bfseries}p{3.0cm}YYY}
\rowcolor{milcVino}\color{white}\bfseries Criterio & \color{white}\bfseries Lo logré & \color{white}\bfseries En proceso & \color{white}\bfseries Necesito apoyo\\[.6mm]
Comprensión del tema & \milcopcion{Explico las ideas con mis palabras.} & \milcopcion{Reconozco las ideas, falta precisión.} & \milcopcion{Necesito repasar lo básico.}\\[.8mm]
\rowcolor{milcGris}Pensamiento computacional & \milcopcion{Usé los pilares al armar mi producto.} & \milcopcion{Usé algunos pilares.} & \milcopcion{Necesito releer la estación 2.}\\[.8mm]
Aplicación y producto & \milcopcion{Mi producto cumple los criterios.} & \milcopcion{Está armado, con ajustes pendientes.} & \milcopcion{Aún está en construcción.}\\[.8mm]
\rowcolor{milcGris}Reflexión 5 dimensiones & \milcopcion{Respondí cada una con honestidad.} & \milcopcion{Respondí algunas con detalle.} & \milcopcion{Mis respuestas son muy generales.}\\[.8mm]
Triángulo de pensamiento & \milcopcion{Conecté las 3 voces con mi trabajo.} & \milcopcion{Conecté al menos una.} & \milcopcion{No las relaciono con mi proceso.}\\[.8mm]
\end{tabularx}}

\vspace{3mm}
\textbf{Hoy entendí que:}\\[1mm]
\linead{\linewidth}\\[1mm]
\linead{\linewidth}

\textbf{Mi compromiso para la próxima sesión es:}\\[1mm]
\linead{\linewidth}\\[1mm]
\linead{\linewidth}

\begin{softbox}{milcMostaza}{milcGris}{Cita de despedida · Marco Aurelio}
\textit{``El obstáculo en el camino se vuelve el camino. Lo que estorba al camino se vuelve camino.''}\\[1mm]
Cada error de hoy, bien observado, se convierte en pieza del camino que pisarás mañana.
\end{softbox}

\small
\begin{tabularx}{\textwidth}{>{\bfseries}p{3.6cm}Y}
\rowcolor{milcGradeDark!12}Nombre completo: & \linead{10cm}\\
Fecha: & \linead{4cm} \quad Curso/grupo: \linead{4cm}\\
Mi firma: & \linead{9cm}\\
\end{tabularx}
\normalsize

\begin{infoband}{milcGradeDark}
\textbf{Para la próxima semana.} Empieza la implementación siguiendo tu cronograma. Mantén la bitácora diaria abierta y registra lo que falla, no solo lo que sale bien. La phronesis del cierre se entrena durante toda la semana, no solo el día del video. Entrega el caso al dueño real al final, antes de presentarlo en clase.
\end{infoband}

\end{document}
